\documentclass[11pt]{article}

% ===================== PACKAGES =====================
\usepackage[a4paper,margin=1in]{geometry}
\usepackage{amsmath,amssymb}
\usepackage{enumitem}
\usepackage{fancyhdr}
\usepackage{xcolor}

% ===================== HEADER & FOOTER =====================
\pagestyle{fancy}
\fancyhf{}
\lhead{CSE 400: Fundamentals of Probability in Computing}
\rhead{Lecture 25 Scribe}
\cfoot{\thepage}

% ===================== TITLE =====================
\title{
\normalsize School of Engineering and Applied Science (SEAS), Ahmedabad University \\
\vspace{0.2cm}
\textbf{CSE 400: Fundamentals of Probability in Computing}\\
\Large Lecture 25: Problem Solving and Tutorial Session (Inequalities)
}
\author{}
\date{}

\begin{document}

\maketitle
\vspace{-2cm}

% ===================== STUDENT INFO =====================
\begin{center}
\begin{tabular}{ll}
\textbf{Name:} & {Khushi Paghadar \hspace{2.5in}} \\
\textbf{Enrollment No:} & {AU2440131 \hspace{2.5in}} \\
\textbf{Email:} & {khushi.p12@ahduni.edu.in \hspace{2.1in}} \\
\end{tabular}
\end{center}

\hrule
\vspace{0.5cm}

% ===================== OBJECTIVE =====================
\section*{Objective}
This scribe serves as an academic reference for exam preparation for Lecture 25

\section*{Overview}
This lecture focuses on applying probabilistic inequalities—primarily Markov, Chebyshev, Jensen, and Chernoff bounds—to solve structured problems. Each solution explicitly follows a step-by-step reasoning process, including identification of the random variable, choice of inequality, and justification.

\vspace{0.5cm}

% ===================== Q1 =====================
\section*{Q1: Markov’s Inequality (Income Problem)}

\textbf{Given:} Random variable $X$ (Annual income), $E[X] = 40{,}000$

\textbf{(a) Find $P(X \ge 120{,}000)$}

Step 1: Identify inequality \\
Only mean is known and $X \ge 0$ \\
$\Rightarrow$ Use Markov’s Inequality:
\[
P(X \ge a) \le \frac{E[X]}{a}
\]

Step 2: Apply
\[
P(X \ge 120{,}000) \le \frac{40{,}000}{120{,}000} = \frac{1}{3}
\]

Final Answer:
\[
P(X \ge 120{,}000) \le \frac{1}{3}
\]

\textbf{(b) Can we bound $P(X < 10{,}000)$}

Reasoning: \\
Markov applies only to $X \ge 0$ and upper tail probabilities. \\
This is a lower-tail probability.

Conclusion: Cannot use Markov’s inequality.

\vspace{0.5cm}

% ===================== Q2 =====================
\section*{Q2: Chebyshev’s Inequality (Resistor Problem)}

Given: $\mu = 100$, $\sigma = 2$, defective outside $[94,106]$

Step 1: Convert interval
\[
|X - 100| \ge 6
\]

\[
k = \frac{6}{2} = 3
\]

Step 2: Apply Chebyshev
\[
P(|X - \mu| \ge k\sigma) \le \frac{1}{k^2}
\]

\[
P(\text{defective}) \le \frac{1}{9}
\]

Step 3: Expected defective out of 10,000
\[
10{,}000 \times \frac{1}{9} \approx 1111
\]

\vspace{0.5cm}

% ===================== Q3 =====================
\section*{Q3: Convexity and Expectation Bound}

Given: $0 \le Y \le 1$, $E[Y] = \mu$

Step 1: Identify tool \\
$g(x) = e^x$ is convex

Step 2: Express $Y$ as mixture
\[
Y = (1-Y)\cdot 0 + Y \cdot 1
\]

Step 3: Apply convexity
\[
e^{tY} \le (1-Y)e^0 + Y e^t
\]
\[
e^{tY} \le (1-Y) + Y e^t
\]

Step 4: Take expectation
\[
E[e^{tY}] \le 1 - \mu + \mu e^t
\]

\vspace{0.5cm}

% ===================== Q4 =====================
\section*{Q4: Jensen’s Inequality (Power Function)}

Given: $X > 0$, $f(x)=x^a$, $a>1$

Step 1: Identify convexity \\
$f(x)$ is convex

Step 2: Apply Jensen
\[
E[X^a] \ge (E[X])^a
\]

\vspace{0.5cm}

% ===================== Q5 =====================
\section*{Q5: Chebyshev General Form}

To show:
\[
P(|X-\mu| \ge k\sigma) \le \frac{1}{k^2}
\]

Derivation:

Apply Markov on $(X-\mu)^2$:
\[
P((X-\mu)^2 \ge k^2\sigma^2) \le \frac{E[(X-\mu)^2]}{k^2\sigma^2}
\]
\[
= \frac{\sigma^2}{k^2\sigma^2} = \frac{1}{k^2}
\]

\vspace{0.5cm}

% ===================== Q6 =====================
\section*{Q6: Production Problem}

(a)
\[
P(X \ge 1000) \le \frac{500}{1000} = \frac{1}{2}
\]

(b)

Step 1:
\[
|X-500| < 100
\]

Step 2:
\[
P(|X-500| \ge 100) \le \frac{1}{100}
\]

Step 3:
\[
P(400 \le X \le 600) \ge 1 - \frac{1}{100} = 0.99
\]

\vspace{0.5cm}

% ===================== Q7 =====================
\section*{Q7: Chernoff Bound Setup}

Given: $X_i = \pm1$ (equal probability)

(a)
\[
E[e^{tX_i}] = \frac{e^t + e^{-t}}{2}
\]

(b) For sum $S_n$
\[
E[e^{tS_n}] = \prod E[e^{tX_i}] = \left(\frac{e^t + e^{-t}}{2}\right)^n
\]

\vspace{0.5cm}

% ===================== Q8 =====================
\section*{Q8: Pairing Problem}

Goal: Bound probability that $\le 30$ mixed pairs

Approach:
Let $X$ = number of mixed pairs

Use Chernoff bound (lower tail)

Reasoning:
Sum of dependent indicators \\
Approximate using Chernoff \\
Bound deviation below mean

\vspace{0.5cm}

% ===================== Q9 =====================
\section*{Q9: Inequality Proof}

Goal:
\[
\sum_{i=1}^{n} \frac{i}{i+1} \ge \frac{n(n+2)}{2(n+1)}
\]

Approach:
Use bounding techniques \\
Recognize monotonic growth

\vspace{0.5cm}

% ===================== Q10 =====================
\section*{Q10: Poisson Tail Bounds}

Given: $X \sim \text{Poisson}(20)$

Want: $P(X \ge 26)$

(a) Markov
\[
P(X \ge 26) \le \frac{20}{26}
\]

(b) Chebyshev
\[
P(X - 20 \ge 6) \le \frac{20}{36}
\]

(c) Chernoff

Step 1: MGF
\[
M_X(t) = e^{20(e^t - 1)}
\]

Step 2:
\[
P(X \ge 26) \le \inf_{t>0} e^{-26t} M_X(t)
\]

Step 3: Optimize
\[
20e^t = 26 \Rightarrow e^t = 1.3
\]

\vspace{0.5cm}

% ===================== TAKEAWAYS =====================
\section*{Key Takeaways}

1. When to use which inequality \\
Markov $\rightarrow$ only mean known, upper tail \\
Chebyshev $\rightarrow$ mean + variance \\
Chernoff $\rightarrow$ tighter bounds \\
Jensen $\rightarrow$ convexity

\vspace{0.3cm}

2. Strategy:
\begin{enumerate}
\item Identify random variable
\item Check given info
\item Select inequality
\item Convert form
\item Apply carefully
\end{enumerate}

\vspace{0.3cm}

3. Exam Tip:

Always justify why a specific inequality is used.

\bigskip
\hrule

\begin{center}
\small \textit{End of Lecture 25 Scribe}
\end{center}

\end{document}